% B05/A023: retained-credit theorem for separated changing cores.
\section*{Fresh rare-switch trees with separated changing cores}

All logarithms are base two.  For a Boolean function $h$ let
\[
 U(h)=\sum_S\widehat h(S)^2
 \log\frac{\sum_{T\supseteq S}\widehat h(T)^2}{\widehat h(S)^2}.
\]

Let $g_0$ and $g_1$ be arbitrary Boolean functions on disjoint physical
variable blocks.  Write
\[
 r_i=\Pr[g_i=-1],\quad \mu_i=1-2r_i,\quad
 q_i=1-\mu_i^2=4r_i(1-r_i),
\]
and
\[
 \theta=\Pr[g_0\ne g_1]=r_0+r_1-2r_0r_1,qquad
 \Delta=r_1-r_0.
\]
The probability here is under independent uniform inputs to the two
disjoint blocks.  Let $z$ be a fresh $d$-bit control block, $p=2^{-d}$,
and define the physical rare switch
\[
 F(z,x,y)=\begin{cases}
 g_1(y),&z=z_\star,\\
 g_0(x),&z\ne z_\star.
 \end{cases}                                                   \tag{1}
\]
Put $\bar r=(1-p)r_0+pr_1$ and
\[
 M(r)=(1-2r)^2\log\frac1{(1-2r)^2}.
\]

\subsection*{Exact two-core identity}

Direct Walsh summation gives
\begin{align}
U(F)={}&M(\bar r)+(1-p)[U(g_0)-M(r_0)]+p[U(g_1)-M(r_1)]\\
&+q_0\left((1-p)^2\log\frac1{1-p}+\frac{pd}{2}-p^2d\right)
 +q_1\frac{pd}{2}\\
&+(2pd-4p^2d)\Delta^2
 +4(p-p^2)\Delta^2\log\frac{\theta}{\Delta^2}.              \tag{2}
\end{align}
The final term is zero when $\Delta=0$.

To verify (2), let $u_T=\widehat g_0(T)$ and
$v_R=\widehat g_1(R)$.  At the empty control frequency, the nonempty
$g_0$-frequencies have coefficient $(1-p)u_T$, the nonempty
$g_1$-frequencies have coefficient $pv_R$, and the empty coefficient is
$1-2\bar r$.  At every nonempty control frequency the corresponding
coefficients are $-pu_T$, $pv_R$, and $2p(r_0-r_1)$.  Because the child
variable blocks are disjoint, their nonempty Fourier supports do not
overlap and no mixed child frequency is present.  Summing the upper masses
over the control supersets gives (2).  Thus (2) is an exact identity, not a
mixture or orthogonality estimate.

The complete signed charge of revealing the $d$ controls, stopping at the
first inactive bit and otherwise reaching $g_1$, is
\[
 J=U(F)-(1-p)U(g_0)-pU(g_1).                                \tag{3}
\]
Every control coordinate is pivotal with probability $2p\theta$, so
\[
 I_{\rm ctl}(F)=2pd\theta.                                  \tag{4}
\]

\subsection*{Uniform node bound}

Set
\[
 L_M=\frac{4(1+2/e)}{\ln2},\qquad
 C_{\rm sep}=3+L_M+\frac2{\ln2}+\frac2{e\ln2}<16.964.       \tag{5}
\]
We prove
\[
 \boxed{J\le C_{\rm sep}I_{\rm ctl}(F).}                   \tag{6}
\]

Let $\phi(\mu)=\mu^2\log(1/\mu^2)$.  Then
$|\phi'|\le L_M/2$.  Since
\[
 |\mu_1-\mu_0|\le1-\mu_0\mu_1=2\theta,
\]
the $M$-Jensen term in (3) is at most $2L_Mp\theta$.
Furthermore
\[
q_0,q_1\le4\theta,qquad \Delta^2\le\theta,qquad
\Delta^2\log\frac{\theta}{\Delta^2}\le\frac{\theta}{e\ln2},\tag{7}
\]
and
\[
(1-p)^2\log\frac1{1-p}\le\frac p{\ln2}.                   \tag{8}
\]
Applying (7)--(8) termwise to (2)--(3) gives
\[
J\le
p\theta\left(2L_M+\frac4{\ln2}+\frac4{e\ln2}\right)
+6pd\theta
\le C_{\rm sep}(2pd\theta),
\]
which proves (6), including all deterministic endpoint cases.

\subsection*{Arbitrary-depth retained-leaf theorem}

Consider any finite rooted tree.  At every internal node $v$, a fresh
$d_v$-bit block implements (1), selecting child $1$ at one word and child
$0$ otherwise.  The complete variable sets of the two child subtrees are
disjoint, and all control blocks are fresh.  Leaves are arbitrary genuine
Boolean functions on mutually disjoint blocks.  Let $w_v$ be the physical
reach probability and $w_\ell$ the leaf probability.

The signed node charges telescope exactly:
\[
\sum_{v\ \mathrm{internal}}w_vJ_v
=U(F_{\rm root})-\sum_{\ell\ \mathrm{leaf}}w_\ell U(g_\ell).\tag{9}
\]
Every selector coordinate is used at exactly one node, and
\[
 I_{\rm sel}(F_{\rm root})
=\sum_{v\ \mathrm{internal}}w_v I_{{\rm ctl},v}.            \tag{10}
\]
Equations (6), (9), and (10) prove the all-depth retained-credit bound
\[
\boxed{
U(F_{\rm root})-\sum_\ell w_\ell U(g_\ell)
\le C_{\rm sep}I_{\rm sel}(F_{\rm root}).}                 \tag{11}
\]
The constant is independent of tree depth, all activation probabilities,
leaf dimensions, leaf spectra, and the number of changing cores.

If the leaves obey $U(g_\ell)\le K I(g_\ell)$, then their physical data
influences also telescope without reuse, yielding
\[
 U(F_{\rm root})\le\max\{C_{\rm sep},K\}I(F_{\rm root}).   \tag{12}
\]

The separation hypothesis is essential to this proof: it deletes every
nonempty Fourier overlap between sibling cores.  Shared variables, reused
address bits, and coherent overlapping sibling spectra produce additional
upper-mass interference not present in (2).  Bounding that overlap defect
by a single-use physical event remains the next universal obligation.
